\section{Conclusion}
\label{sec:conclusion}

We introduced \ours, a graduated adversarial stress-test harness that turns discarded adversarial null results into noise- and validity-calibrated robustness boundaries, applied identically across five model families, four failure modes, and six intensity levels (2{,}879 real API responses). Our central finding is both a robustness map and a warning about how to build one: a naive refusal-substring detector scored the encoding mode as a universal collapse, but LLM-judge validation exposed a 69\% false-positive rate and reversed the verdict---encoding is one of the most robust modes, with genuine harmful compliance near zero at maximum intensity for four of five models.

\para{Key findings.} The validated map ranks models \gpt $>$ \gemma $>$ \llama $>$ \qwen $>$ \mistral, shows robustness is multi-dimensional (encoding is orthogonal to a correlated jailbreak--persona--bias cluster), identifies 11 of 20 (model, mode) invariants, and finds preference-optimized and deliberative models more robust than fine-tuning-aligned ones (0.771 vs.\ 0.581). The take-away is a single sentence: adversarial null results are quantifiable robustness boundaries, not experimental weaknesses---\emph{provided the behavioral metric is itself validated.}

\para{Future work.} Three directions follow. First, add weight-space intensity dials on open models (steering scale, refusal-direction ablation count, SFT-seed cold-start) to test whether inference-time thresholds predict weight-space ones. Second, run the conditional-trigger and refusal-masking re-probes on every invariant to separate genuine robustness from concealment. Third, scale to $\geq 15$ models to power the alignment-method contrast and correlate thresholds with mechanistic predictors such as safety-layer band and misalignment-direction magnitude.
